\documentclass[11pt]{amsart}
\usepackage[margin=1in]{geometry}
\usepackage{amsmath,amssymb,array,booktabs,longtable,xurl}
\usepackage[hidelinks,hypertexnames=false]{hyperref}
\newcommand{\Z}{\mathbb Z}
\newcommand{\N}{\mathbb N}
\newcommand{\D}{\mathrm{D12}}
\newcommand{\record}[1]{\nolinkurl{#1}}
\newcommand{\claimid}[1]{}
\newtheorem{theorem}{Theorem}
\newtheorem{lemma}[theorem]{Lemma}
\newtheorem{proposition}[theorem]{Proposition}
\newtheorem{corollary}[theorem]{Corollary}
\theoremstyle{definition}\newtheorem{definition}[theorem]{Definition}
\setlength{\emergencystretch}{3em}
\renewcommand{\qedsymbol}{\rule{0.6em}{0.6em}}
\renewcommand{\arraystretch}{1.12}
\raggedbottom
\title{Supplementary proof notes for the D12 section of article 6.2}
\author{Tatsuhiko Takimoto}
\address{Independent researcher}
\date{Draft v2, 21 October 2026. Unpublished.}
\begin{document}
\begin{abstract}
We study literal linear generalized Stone--Wales paths in the
positive-parameter D12 family. Two separately derived proof routes give
the criterion $r\mid(s-1)$ or $r\mid(s+1)$ or $|r-3s|=1$.
Route A uses a bespoke symbolic executor; Route B uses explicit
lattice strips, rational infeasibility certificates and induction on
an arbitrary number of strips. We state the geometric, code and prose
dependencies of each route. Finite graph audits compare the symbolic
decisions with an eager graph construction and a definition-based path
search. Both universal arguments use prose geometry and induction;
the ray consequence follows by elementary arithmetic. The manuscript
is an unpublished draft, with human proof peer review and priority
comparison open.
\end{abstract}
\maketitle

\section{Introduction}
Graver's catalogue supplies the D12 construction and its parameters
\cite{Graver}. Bille, Buchstaber and Spodarev formulate the graph-path
conditions and the broader fullerene questions used here \cite{BBS}.
The standard cyclic edge-connectivity fact recalled by Kutnar and
Maru\v{s}i\v{c} supplies the local continuation premise \cite{Connectivity}.
Our proposed contribution concerns an exact path criterion for this
existing catalogue family and the finite computations supporting its
symbolic proof route. The family construction and the definition are
inherited. Within the literature searched so far, no matching complete D12
criterion has been identified. This bounded search does not establish
priority; the comparison remains open.

\clearpage
\section{What is proved, and how}\label{sec:guarantees}
\begin{center}\small
\begin{tabular}{p{.22\textwidth}p{.28\textwidth}p{.38\textwidth}}
\toprule Claim & Present proof basis & Scope\\\midrule
Theorem~\ref{thm:d12-classification} & Two prose/computer-assisted routes & Positive integer $r,s$ in the stated quotient convention, with the model and soundness premises below.\\
Route A & Exact symbolic region splitting and cylinder induction & All ten representative starts; six semantic obligations in Section~\ref{sec:trust}.\\
Route B & Explicit strips, rational infeasibility certificates and prose induction & Universal necessity and sufficiency; global simplicity covers arbitrary strip counts.\\
Finite audits & Eager and lazy graph constructors, literal-path searches and certificate checks & The input ranges and overlapping counts stated below.\\
Ray consequence & Conditional prose arithmetic & Every eventually positive integer affine ray, assuming the classification.\\
Auxiliary start counts & Empirical prediction & The observed inputs; the stated universal boundary count is open.\\
\bottomrule
\end{tabular}
\end{center}
The two routes are separately derived on the same positive-parameter
domain, without sharing their D12 constructions, proofs or code.
Their shared model and semantic premises are stated below; human proof
peer review, Lean formalization and the priority comparison remain open.
The supplementary \record{VERIFICATION-RECORD.md} contains the roles,
outcome-exposure qualification, reproduction provenance and history.

This third manuscript follows an A3 classification manuscript and a
ray-periodicity manuscript \cite{A3,Lines}. The latter includes D12 in
its catalogue-ray corollary but deliberately excludes a full D12
classification. The present criterion would close that D12 subproblem;
it does not change the second manuscript's general geometric theorem.
K3, C1, O2 and D14 have Lean graph results. Those
results provide formalization patterns, not a Lean proof of D12.

\section{The model and the literal path}
\begin{definition}[D12 quotient convention]\label{def:d12-model}
Let $r,s\in\Z_{\ge1}$ and $L=r+2s$. In the triangular lattice with
unit differences
\[
 (1,0),(0,1),(-1,1),(-1,0),(0,-1),(1,-1),
\]
take two chart-tagged copies of the closed hexagon
\[
 H_L=\{(x,y)\in\Z^2:-L\le x\le L,\ 0\le y\le2L,
                         \ 0\le x+y\le2L\}.
\]
Its counterclockwise vertices are
$(0,0),(L,0),(L,L),(0,2L),(-L,2L),(-L,L)$.
Mark each side at distance $r+s$ from its initial vertex. The twelve
boundary points, alternating vertices and marks, receive labels
\[
 (11,0,1,2,\ldots,10)\quad\hbox{and}\quad(0,11,10,\ldots,1)
\]
in the two charts. Glue equally labelled segments in reverse order,
preserving lattice distance along the segment. Vertices are equivalence
classes of chart sites; adjacency is a unit edge within a common chart.
Write $G(r,s)$ for the resulting simple graph.
\end{definition}
\claimid{def:d12-model}

This normalized D12 signature is obtained by translating
the first chart by $(r+s,0)$ and cyclically shifting its corner list
so that both underlying hexagons equal $H_L$. All 48 coordinate-form
identities and the labels agree by exact coefficient comparison.
The catalogue convention has coordinate lengths $r+s$ and $s$ and
carbon order $12(r+2s)^2$ \cite[p.~12]{Graver}.

\paragraph{Geometric premise.}\label{lem:geometry}
The intended quotient is a simple fullerene dual on a sphere, with
$6L^2+2$ vertices, twelve degree-five vertices at the marked labels,
and degree six elsewhere. Pairing the entire boundaries of two disks
in opposite order explains the sphere and the absence of extra
within-chart identifications. Each marked label joins a hexagon angle
$2\pi/3$ to a straight boundary angle $\pi$. All other sites have total
angle $2\pi$. These geometric model premises are supported by the
finite model audit and belong to the computer-assisted argument's
prose surface correspondence.
\claimid{lem:geometry}

A fullerene is cyclically 5-edge-connected, as recalled in
\cite{Connectivity}. A nonfacial dual triangle would give a cyclic
three-edge cut in the primal graph; hence every dual triangle is facial.
An edge of the simple triangulated sphere therefore has exactly its two
facial common neighbours. This external connectivity fact supplies the local continuation premise.

\begin{definition}[Literal linear gSW path]\label{def:literal}
A literal path is a sequence $p_0,\ldots,p_{n-1}$ with even $n\ge4$,
all vertices distinct, with consecutive and distance-two vertices
adjacent, and
\[
 \deg p_0=\deg p_{n-1}=5,\qquad
 \deg p_1=\deg p_{n-2}=6.
\]
\end{definition}
\claimid{def:literal}
This is the path condition of \cite[Definition 1(a)]{BBS}, using a
vertex count rather than an edge count. Even vertex count is part of
the definition. Internal degree-five vertices are allowed. We classify
this graph predicate; no chemical activation barrier or experimental
isomerization conclusion is inferred.

\begin{lemma}[Forced prefixes]\label{lem:forced}
Suppose every edge of a finite graph has exactly two common neighbours.
After a valid oriented starting triangle $(p_0,p_1,p_2)$ is fixed,
every literal path extending it is a prefix of the continuation
\[
 p_{i+1}\in(N(p_{i-1})\cap N(p_i))\setminus\{p_{i-2}\}.
\]
The right side has one element. A path exists from that start precisely
when a legal endpoint is reached strictly before the first repeated
vertex. A legal endpoint has odd index at least three, degree five,
and a degree-six predecessor.
\end{lemma}
\claimid{lem:forced}
\begin{proof}
Skip adjacency and distinctness force the other common neighbour.
A legal prefix satisfies every condition in Definition~\ref{def:literal}.
A first repeated vertex prevents all longer prefixes from being simple.
At a tie, repetition invalidates the prefix. Since the graph is finite,
a first repetition or earlier legal endpoint is eventually reached.
\end{proof}

\section{Classification and route A: the finite region cover}
\begin{theorem}[Classification with disclosed proof dependencies]\label{thm:d12-classification}
For the positive-parameter quotient convention above, the two received
proof routes assert that $G(r,s)$ has a literal path if and only if
\[
 \boxed{\quad
 r=1\quad\hbox{or}\quad
 (r=2\ \hbox{and }s\hbox{ odd})\quad\hbox{or}\quad
 \bigl(r\ge3\ \hbox{and }(s\equiv\pm1\pmod r\ \hbox{or }|r-3s|=1)\bigr).
 \quad}
\]
Equivalently, for all positive $r,s$, the criterion is
$r\mid(s-1)$ or $r\mid(s+1)$ or $|r-3s|=1$. Indeed, divisibility by
one is automatic; modulo two, $+1$ and $-1$ both mean odd $s$, while
$|2-3s|=1$ adds only $s=1$; for $r\ge3$ the formulas are identical.
\end{theorem}
\claimid{thm:d12-classification}
Route A uses the six semantic obligations in Section~\ref{sec:trust};
route B uses the explicit construction, rational certificates and prose
induction in Section~\ref{sec:route_b}. These are alternative routes on
the shared model, not a combined list of premises. The domain is the
positive-parameter quotient convention.

\paragraph{Symmetry reduction.}\label{lem:symmetry}
The label permutations $u\mapsto u+2$ and $u\mapsto1-u$ are realized
by affine lattice isometries, with the second interchanging the charts.
The symbolic corner check verifies the twelve corner identities in
each chart. The induced group is transitive on the twelve degree-five
labels. It suffices to consider the ten candidate oriented triangles
at label zero, numbered as in \record{start_triples()}. A candidate
with degree-five second vertex is invalid. There are 120 valid starts
when $s>1$ and 96 when $s=1$ in the supplied model.
\claimid{lem:symmetry}

\paragraph{Finite regions.}\label{prop:regions}
On each region the executor splits integer affine sign conditions until
an invalid start or the earlier of a legal endpoint and repetition has
a fixed symbolic description. The symbolic core summaries give:
\claimid{prop:regions}
\begin{center}\small
\begin{tabular}{p{.25\linewidth}r p{.57\linewidth}}
\toprule
Region & Leaves & Positive canonical starts and conditions\\\midrule
$s\ge2,\ r\ge s+2$ &61& Start 0 at $r=3s-1$; start 3 at $r=3s+1$.\\
$s=1,\ r\ge3$ &20& Starts 2,4,6,8; additionally start 3 at $r=4$.\\
$s=r-1,\ r\ge3$ &18& Starts 1,3,5,7.\\
$s=r,\ r\ge3$ &10& None.\\
$s=r+1,\ r\ge3$ &12& Starts 0,2,4,8.\\
$s\ge r+2,\ r\ge3$ &6& None for starts 1,2,5,6,7,9.\\
$W_1,\ldots,W_9$ &304& Remaining starts 0,3,4,8; positive at the two modular boundaries described below.\\\bottomrule
\end{tabular}
\end{center}
Here $W_q=\{r\ge3:qr+2\le s\le(q+1)r+1\}$. These nine wedges cover
$r+2\le s\le10r+1$ for the four remaining starts. The two exceptional
lines in the first row have endpoint indices $56s-11$ and $56s+11$.
The count 304 is a count of symbolic leaves, not of fullerene isomers.
The leaves are not assumed disjoint across alternative starting triangles.

For $s\ge10r$, the remaining four starts use the cylinder argument in
Section~\ref{sec:cylinders}. There is an overlap with $W_9$, rather than
a gap. For $r=2$, the non-cylinder starts are treated symbolically and
the cylinder calculation applies for $s\ge20$; the finite complement is
$s\le19$. For $r=1,s\ge2$, canonical start 1 is positive in the boundary
certificate. The point $(1,1)$ is decided directly. These regions cover
all positive integer parameters. Together with the terminal events and
symmetry reduction, they give the stated classification under the
soundness obligations.

\section{Unbounded passage counts and cylinder exits}\label{sec:cylinders}
\paragraph{Cylinder state.}\label{lem:cylinders}
The passage count is not bounded uniformly in $s/r$. Instead a state
family $S(a)$ recurs after four runs as $S(a+2r)$. The common state shape
for starts 0 and 3 in chart 0, and start 4 in chart 1, is
\[
 ((a+1,-1),(a,0),(a+1,0)).
\]
Start 8 in chart 1 uses $((a+1,-1),(a+1,0),(a,0))$.
The index increment is $12r+20s$ for starts 0,3,4 and $8r+20s$ for
start 8. The transient has no earlier legal endpoint or repetition;
the lap calculation returns to the same state family and is cone-free.
\claimid{lem:cylinders}

\begin{center}\small
\begin{tabular}{c c c c c}
\toprule
Start&$a_0$&Last lap bound $h$&Positive exit planes&Residue\\\midrule
0&$2r$&$s-2r-2$&$s-2r-1,\ s-r-1$&$s\equiv1$\\
3&$2r-1$&$s-2r-2$&$s-2r,\ s-r$&$s\equiv-1$\\
4&$3r+s$&$2s-r-2$&$2s-r-1,\ 2s-1$&$s\equiv1$\\
8&$2r$&$s-r-2$&$s-r-1,\ s-1$&$s\equiv1$\\\bottomrule
\end{tabular}
\end{center}
Residues in the last column are modulo $r$.

\paragraph{All previous laps.}\label{lem:history}
A negative exit leaf exhibits a repetition in the transient/last-two-lap
history before any legal endpoint. For positivity that bounded history
is insufficient: every earlier lap must be excluded. The four-variable
calculation compares the exit against a lap with arbitrary integer
parameter $b\in[a_0,a-2r]$, the transient against every such lap,
and two laps separated by at least $2r$. It also checks transient/exit
and internal coincidences. Relaxing the actual arithmetic progression
to every integer $b$ gives a sufficient overapproximation for collision
exclusion. A proof of the implementation of these comparisons remains
one of the universal obligations.
\claimid{lem:history}

\begin{lemma}[Exit arithmetic]\label{lem:exit-congruence}
On $r\ge2,s\ge10r$, accept the geometric cylinder and terminal-leaf
premises of the table. Then its exit conditions reduce respectively
to $s\equiv1,-1,1,1\pmod r$.
\end{lemma}
\claimid{lem:exit-congruence}
\begin{proof}
Let $a_j=a_0+2rj$. There is a unique first value $a_f>h$, and
$h+1\le a_f\le h+2r$. The two positive planes in each row lie in this
window, differ by $r$, and represent the two lifts of one residue
modulo $r$ to residues modulo $2r$. Thus $a_f$ equals one of the two
planes exactly when that plane is congruent to $a_0$ modulo $r$.
Substituting the four rows gives the stated residues. The bound
$s\ge10r$ ensures $a_f\ge a_0+4r$, so the two prior laps used by the
exit calculation exist.

These inequalities and the exact reduction of the positive-leaf
constraints to the displayed planes were checked independently of
the geometric executor. Substitute each plane into all its constraints,
then write $r=2+u$, $s=10r+v$, $u,v\ge0$. Every required expression has
nonnegative constant and coefficients. In the other direction, each
positive leaf explicitly contains both $a-c\ge0$ and $c-a\ge0$ for
its plane $a=c$, including the window constraint where needed.
This proves the arithmetic identities; it does not prove the
geometric premise that the listed terminal leaves are correct.
\end{proof}

\section{Finite checks and input ranges}
The exact-rational symbolic recipe has thirteen steps and seven
certificate summaries. It covers thirty mirror graphs, the symmetry
identities, finite-region and wedge summaries, cylinder and full-history
positivity summaries, 452 cylinder sample/start checks and exit
arithmetic. The comparison checks JSON structures after removing runtime
fields. Wedges above $q=9$ are exploratory inputs separate from the
finite region cover; the arbitrary tail uses the cylinder argument.

\paragraph{Separate finite implementation.}\label{prop:independent-finite}
A separately written eager constructor enumerates every lattice site
of the two normalized hexagons, forms the boundary quotient by
union--find, and builds all unit edges. It then searches all simple
literal prefixes by backtracking common neighbours; its search routine
does not assume unique continuation. The pilot $r,s\le3$ is included
in the subsequent $r,s\le6$ grid of 36 pairs and 4,176 starts.
The finite complement and known cutover tests add $r=2,1\le s\le21$
and $r=3,s\in\{29,30,31,32\}$. After removing overlap there are
55 parameter pairs and 6,456 starts: 33 positive pairs, 22 negative
pairs, and 1,800 successful oriented starts. Every vertex partition,
edge set, start set and start outcome agrees with the lazy decider.
The largest checked graph has 26,936 dual vertices.
\claimid{prop:independent-finite}

At the same 36 grid points the two-, three- and four-variable executors
were specialized to exact point regions. All 1,080 canonical-start
comparisons agree with the eager graph on invalidity, first event
kind and its exact index. These are overlapping finite checks of known
inputs. Constructor, search and symbolic-event comparisons have
these finite ranges.

\paragraph{Discovery and fixed-prediction ranges.}\label{prop:holdout}
The 1,600 discovery points are $r,s\le40$. The saved
295 test points all have $\max(r,s)\ge41$; 74 are positive. The saved binary decisions and successful-start counts agree
with the fixed predictions.
\claimid{prop:holdout}

The generated test set has nine pairs beyond its stated
$r,s\le400$ range, with maxima $r=441$, $s=453$. These are the actual
input limits of the saved comparison.

The wedge prediction fits event formulas at $q=1,\ldots,6$
and tests them at $q=7,\ldots,24$. Their saved terminal-event formulas
agree with the fixed predictions; the arbitrary tail uses the cylinder
argument. The planned range $q=25,\ldots,30$ has no executed tests.

\paragraph{Auxiliary counts.}\label{conj:start-counts}
The saved prediction also gives successful-start counts, agreeing on
all discovery and test records and on the independent finite tests.
In particular it predicts 96 successful starts for $r=1,s\ge3$.
That universal boundary count remains an auxiliary prediction.
We retain these counts as an auxiliary prediction, not as part of
Theorem~\ref{thm:d12-classification}. Counts refer to starting oriented
triangles, not to distinct paths or fullerene isomers.
\claimid{conj:start-counts}

\section{Route B: explicit strips and rational certificates}\label{sec:route_b}
Route B \cite{ProofB} proves necessity and sufficiency separately using explicit
strip formulas, seam maps and rational certificate inputs.

\paragraph{Coordinates and starts.}
Put $D=r+s$, $B=r+3s$, $C=2r+3s$ and retain $L=r+2s$.
B's chart zero is
\[
 -C\le x\le s,\qquad 0\le y\le2L,\qquad -D\le x+y\le B;
\]
chart one is $H_L$. Translating chart zero by $(D,0)$ gives precisely
the two hexagons and labelled boundary convention of
Definition~\ref{def:d12-model}. Thus this is a coordinate change, not a
different parameter convention. Passing from affine coordinate and
label identities to the spherical quotient remains a prose argument.
With $T(x,y)=(-y,x+y)$, the symmetries are
\[
 R:\ (T+(s,D),\ T^5+(-L,L)),\qquad
 J:\ (T^5+(-s,s),\ T+(s,0)),
\]
where $J$ exchanges charts; their label actions are $i\mapsto i+2$
and $i\mapsto1-i$. Let $P=R^3$ and $K=R^{-1}J$.
The arithmetic checks verify 48 corner correspondences for $R,J$ and
24 each for $P,K$. Transitivity reduces necessity to ten starts at
label zero. \emph{B's start numbering is local to B}: its residual
starts are $0,3,4,6$, whereas A uses $0,3,4,8$ for its cylinder domain.
No equality of start numbers across implementations is assumed.

\paragraph{Necessity: short strips.}\label{prop:route_b-necessity}
Write a strip as $(f,A,d,e,n)$, representing vertices
\[
 A+\lfloor j/2\rfloor d+(j\bmod2)e,\qquad 0\le j<n,
\]
in chart $f$; $n$ counts vertices. Unit differences $d,e,d-e$ give
triangles with recurrence $p_{j+3}=p_{j+1}+p_{j+2}-p_j$.
For negative candidates $r,s\ge2$, B covers the following regions:
\begin{center}\small
\begin{tabular}{ll}\toprule
Region & Excluded B-starts\\\midrule
$r=s$ & all ten\\
$s+2\le r\le3s-2$ & all ten\\
$r=3s$ & all ten\\
$r\ge3s+2$ & all ten\\
$s\ge r+2$ & $1,2,5,7,8,9$\\\bottomrule
\end{tabular}
\end{center}
There are 29 explicit prefix candidates. Each certifies chart
containment, seam continuation, a repeated vertex and exclusion of
every earlier possible odd-index degree-five endpoint. The repeated
vertex need not be the globally first repeat: no earlier legal endpoint
and any subsequent repetition suffice for nonexistence. The 46
region--start cover records and 2,464 odd-vertex exclusions make this
a parameter-domain argument, not a sample fit.

For each infeasible linear system $f_i(r,s)\ge0$, nonnegative rational
weights satisfy $\sum_iw_if_i=-c$ with $c>0$ and all variable
coefficients zero. Adding the coefficients verifies infeasibility
without trusting the search for the weights. Equalities are represented
by two inequalities. Integer coverage splits on the first violated
facet, using $\neg(f\ge0)\iff -f-1\ge0$ for integer-valued $f$.

\paragraph{Necessity: arbitrarily many strips.}
For B-starts 3 and 4, put $\varepsilon=-1,+1$, respectively.
The repeated pair is $P^jE_{jr},P^jO_{jr}$, where
\begin{center}\small
\begin{tabular}{ccllll}\toprule
Strip & $f$ & $A$ & $d$ & $e$ for $-1;+1$ & $n$\\\midrule
$E_h$ &0&$(h,0)$&$(0,1)$&$(-1,1);(1,0)$&$2(B-h)$\\
$O_h$ &1&$(-L,r+4s-h)$&$(1,0)$&$(0,1);(1,-1)$&$2(L+h+r)$\\\bottomrule
\end{tabular}
\end{center}
The one-pair identities hold for $h\ge0$, $r,s\ge2$, with
$s-h\ge r$ for $-1$ and $s-h\ge r+2$ for $+1$.
They prove containment, seam transitions and absence of odd-index
degree-five points. Because $P$ is an automorphism, induction gives
the sequence for arbitrary $j$, not just tested lap counts.

In the remaining negative region $s\ge r+2$, with
$s\not\equiv\pm1\pmod r$, choose
\[
 m=\begin{cases}\lfloor s/r\rfloor&\varepsilon=-1,\\
                  \lfloor(s-2)/r\rfloor&\varepsilon=+1,
   \end{cases}\qquad h=mr,\quad v=s-h.
\]
Here $m\ge1$; $v$ is either in $[2,r-2]$, or is $0$ for $-1$
and $r$ for $+1$. In the orientation pulled back by $P^m$, the final
pair reaches a terminal strip beginning at $(-C,r+5s-h)$.
Its last point equals a previously visited point of $PE_{h-r}$,
at local index $4r-2$ for $-1$ and $4r-1$ for $+1$.
The bound $h\le s\le3s$ puts that index inside the earlier strip.
All preceding odd-index degree-five candidates are excluded by the
linear certificates. B-starts 0 and 6 have even initial prefixes of
length $2D$ and $2L$, followed by $K$- and $R$-images of start 4.
Any earlier collision with those prefixes only strengthens negativity.
This closes all ten starts; the positive boundary cases excluded from
the negative domains are supplied next.
\claimid{prop:route_b-necessity}

\paragraph{Sufficiency and global simplicity.}\label{prop:route_b-sufficiency}
B supplies the following positive witnesses, with vertex counts:
\begin{center}\small
\begin{tabular}{lll}\toprule
Parameters & Shape & Vertices\\\midrule
$r=1$, $s\ge1$ & one strip & $6s+4$\\
$s=1$, $r\ge2$ & two strips & $4r+6$\\
$r=3s+1$, $s\ge1$ & four strips & $56s+12$\\
$r=3s-1$, $s\ge2$ & four strips & $56s-10$\\
$s=qr\pm1$, $r\ge3$, $q\ge1$ & alternating strips &
 $2q(3r+5s)+2$ for $+1$, $2q(3r+5s)$ for $-1$\\\bottomrule
\end{tabular}
\end{center}
The $+1$ construction extends to $r=2$ by 552 coefficient checks
and the accompanying simplicity argument; $q=0$ is the $s=1$ case.
The $r=1$ strip has even vertices $(-j,j)$ and odd vertices
$(-j-1,j)$, which have distinct coordinate sums. Its endpoint is
label 8 and its adjacent vertices have degree six.

For the four boundary families, exact one-variable inequalities check
containment, endpoints and intersections of every pair of even/odd
rows, including identified seams. Nonparallel directions have
determinant $\pm1$; parallel directions reduce to interval overlap.
For the unbounded congruence family, the prose proof separates
columns and rows by $2r$, separates the two parities into disjoint
chart ranges, and assigns ownership of seam vertices to avoid counting
a boundary point in both blocks. These are global simplicity
arguments, not an assertion based only on the final two strips.
Explicit unit triangles prove existence without the external
cyclic-connectivity theorem. Internal degree-five sites are allowed.
\claimid{prop:route_b-sufficiency}

\paragraph{Rational and finite audits.}\label{prop:route_b-reproduction}
The finite audit contains 120 graphs, 14,160 starts and 1,180 canonical
coordinate traces, with 57 positive and 63 negative pairs, overlapping
the previous graph audit. A separate arithmetic implementation checks
9,542 rational certificates, 1,094 one-variable integer conditions
and 552 coefficients; it rederives 29 short geometries, 2,464 exclusions
and 46 covers. The long-strip argument uses the displayed induction
and primary checker.
\claimid{prop:route_b-reproduction}

\paragraph{Derivation scope.}\label{claim:independence}
The derivation and outcome-exposure qualifications are documented in
the supplementary record; the universal route uses the explicit
geometric and arithmetic premises stated here.
\claimid{claim:independence}

\section{Semantic premises for the two routes}\label{sec:trust}
The following six obligations delimit route A of the universal
classification. Route B has its own strip and certificate semantics.
\begin{enumerate}
\item The quotient convention must faithfully represent the catalogue
family, and its all-parameter simple-sphere and degree properties must
hold. The geometric explanation, exact chart normalization and finite
checks are distinguished from a formal proof of this correspondence.
\item Affine passage rows must enumerate exactly the forced vertices,
with correct seam transfers, endpoint parity and predecessor degree.
\item Polyhedral sign tests must be sound and integer branch splits
exhaustive. The code uses exact rational vertices and recession rays.
The emptiness shortcut assumes pointed polyhedra; callers must preserve
that property when adding the auxiliary variables and constraints.
\item Coincidence tests must cover same-chart, seam and parallel
intersections and return the earliest repeated index. A tie must be
resolved in favour of repetition.
\item Cylinder transients, lap induction, exit coverage and arbitrary
previous-lap collision exclusion must connect to the actual walk.
\item The finite parameter-region cover and the symmetry reduction
must cover every valid start, including the small boundary cases.
\end{enumerate}
The symbolic summaries contain branch constraints, terminal indices,
flags and sample readbacks. The universal route uses the six semantic
premises above to connect them to the actual graph.

For route B, the corresponding dependencies are the all-parameter
surface construction and catalogue correspondence; the quotient
automorphisms and exhaustive start reduction; the external cyclic
5-edge-connectivity theorem for necessity; the strip formulas and
unbounded induction; and the global simplicity and seam ownership of
positive witnesses. Rational weighted sums can be checked independently
of their search procedure, but relating their input inequalities to the
actual graph remains a geometric and semantic obligation. The separate
arithmetic audit rederives the short-strip geometry, not the entire
long-strip geometry. B's explicit positive triangles do not require
the external connectivity theorem.

\section{Relation to affine-ray periodicity}\label{cor:d12-rays}
\begin{corollary}[Conditional arithmetic consequence]
Assume Theorem~\ref{thm:d12-classification}. On every eventually positive
integer affine ray $r(t)=at+b$, $s(t)=ct+d$, path existence is eventually
periodic. If $a>0$, the decision is eventually constant. If $r=m$ is
constant, a valid period divides $m/\gcd(m,c)$ for $m\ge3$; the cases
$m=1,2$ have periods dividing one and two, respectively.
\end{corollary}
\claimid{cor:d12-rays}
\begin{proof}
Eventual positivity forces $a,c\ge0$. For $a>0$, set
$\Delta_\varepsilon=a(d-\varepsilon)-bc$, $\varepsilon\in\{-1,1\}$.
If $r(t)\mid s(t)-\varepsilon$, then $r(t)\mid\Delta_\varepsilon$.
A nonzero fixed $\Delta_\varepsilon$ excludes this once
$r(t)>|\Delta_\varepsilon|$. If it is zero, the ratio
$(s(t)-\varepsilon)/r(t)=c/a$ is constant, so divisibility holds
exactly when $a\mid c$. The condition $|r(t)-3s(t)|=1$ is eventually
constant as well: its affine coefficient either vanishes or its
absolute value eventually exceeds one. The boundary $r\le2$ eventually
disappears. Thus the decision stabilizes.

For constant $r=m\ge3$, the two congruences are periodic with period
$m/\gcd(m,c)$. The exceptional affine equality is constant or eventually
absent. For $m=1$ the criterion is always true, and for $m=2$ it is the
parity of $s$. These arguments give elementary threshold bounds from
the coefficients. They are prose arithmetic, not a certified geometric
compiler or a least-period computation.
\end{proof}
The second manuscript already supplies a broader D12 ray statement.
Here the conditional closed classification explains its decisions
arithmetically; no new general periodicity principle is claimed.

\section{Further mathematical obligations}
A direct formalization can use a quotient graph and path-transport
interfaces with generated strip expressions and kernel-checked algebra.
The central obligations are graph semantics, exhaustive starts and
arbitrary-length strip induction. Positive congruence witnesses require
collision exclusion against every previous strip and across seams;
a fixed last-two-strips check is insufficient. Rational witnesses
provide arithmetic input while their relation to surface geometry
requires proof.

\section*{Authorship and AI assistance}
The author used AI and AI system for mathematical exploration,
proof drafting, implementation, finite and arithmetic checks,
reproduction, evidence organization and English exposition. The author
is responsible for the manuscript and its stated verification scope.
AI systems are not listed as authors.

\begin{thebibliography}{9}
\bibitem{Graver} J. E. Graver,
\emph{Catalog of All Fullerenes with Ten or More Symmetries},
DIMACS Series 69 (2005), 167--188; author version with corrections,
D12 entry on PDF page 12.
\bibitem{BBS} A. Bille, V. Buchstaber, and E. Spodarev,
\emph{Some open mathematical problems on fullerenes},
arXiv:2410.19322v1 (2024), Definition 1(a).
Published version (2025), doi:10.1021/acs.jcim.4c01997.
\bibitem{Connectivity} K. Kutnar and D. Maru\v{s}i\v{c},
\emph{On Cyclic Edge-Connectivity of Fullerenes},
arXiv:math/0702511v1 (2007). The cited cyclic 5-edge-connectivity fact
is recalled there as known.
\bibitem{A3} T. Takimoto,
\emph{A two-stage classification of linear generalized Stone--Wales paths
in Goldberg--Coxeter images of the dodecahedron}, version 8.1 (2026).
\bibitem{Lines} T. Takimoto,
\emph{Eventual periodicity of linear generalized Stone--Wales paths in
affine fullerene families with convex lattice-face decompositions},
version 5.0 (2026), Corollary 1b and the D12 partial-result discussion.
\bibitem{ProofB} Unpublished research note,
\emph{D12 classification: necessity and sufficiency with rational
certificates}, 20 October 2026. Exact provenance and authorship roles
are supplied in the supplementary record.
\end{thebibliography}
\end{document}
